\chapter{القياس والزمن}\label{ch:g1:measure}

أيّ القلمين أطول؟ وأيّ يوم اليوم؟ وكم تكلّف الحلوى؟ في هذا الفصل
نقارن الأطوال ونقيسها بالسنتيمتر، ونرتّب أيام الأسبوع، ونقرأ
الساعات الكاملة على الساعة، ونعدّ النقود
الصغيرة.

\section{أطول وأقصر}

\begin{method}[مقارنة الأطوال]\label{met:g1:measure:compare}
لمقارنة طولين، صُفَّهما \emph{من طرف واحد مشترك}:
\begin{enumerate}
  \item ضع الشيئين جنبًا إلى جنب، بادئين من الخط نفسه؛
  \item فالذي يمتد أبعد هو \emph{الأطول}، والآخر هو
        \emph{الأقصر}؛
  \item وإن بلغا الموضع نفسه فهما \emph{متساويان في
        الطول}.
\end{enumerate}
ومطابقة البداية هي السرّ كله --- فالمقارنة من بدايات ملتوية تخدع
العين.
\end{method}

\begin{omfigure}
\begin{tikzpicture}[scale=0.9]
  \draw[very thick, omDef] (0,0.6) -- (3.4,0.6);
  \draw[very thick, omThm] (0,0) -- (2.2,0);
  \draw[gray, dashed] (0,-0.2) -- (0,0.8);
  \node[right, font=\small] at (3.6,0.6) {الأطول};
  \node[right, font=\small] at (3.6,0) {الأقصر};
\end{tikzpicture}

{\small عودان مطابقان عند اليسار: الأزرق يمتد أبعد،
فهو الأطول.}
\end{omfigure}

\begin{definition}[السنتيمتر]\label{def:g1:measure:cm}
لنقول \emph{كم} يبلغ الطول، نقيس بمسطرة مدرَّجة
\emph{بالسنتيمتر}\index{سنتيمتر} (سم). ضع العلامة $0$ من المسطرة
عند بداية الشيء، واقرأ العدد عند نهايته.
\end{definition}

\begin{omfigure}
\begin{tikzpicture}[scale=1.0]
  % ruler body extends a little past the last mark so no numeral is cut
  \draw[thick] (-0.4,0) rectangle (8.4,0.7);
  \foreach \i in {0,...,8} {
    \draw[thick] (\i,0.7) -- (\i,0.4);
    \node[font=\footnotesize] at (\i,0.2) {\i};
  }
  \draw[very thick, omThm] (0,1.15) -- (5,1.15);
  \node[omThm, above, font=\small] at (2.5,1.22) {القلم: $5$ سم};
\end{tikzpicture}

{\small القياس: العلامة $0$ عند البداية، ثم اقرأ عند النهاية. طول
هذا القلم $5$ سم. (والمليمتر والمتر يأتيان في
\cref{ch:g3:measure}.)}
\end{omfigure}

\section{أيام الأسبوع}

\begin{example}[الأيام السبعة]\label{ex:g1:measure:days}
في الأسبوع $7$ أيام، وهي دائمًا بالترتيب نفسه:
\[
\text{الاثنين، الثلاثاء، الأربعاء، الخميس، الجمعة، السبت،
الأحد.}
\]
وبعد الأحد يبدأ الاثنين من جديد. و\emph{أمس} هو اليوم السابق،
و\emph{غدًا} هو اليوم التالي: فإن كان اليوم هو الأربعاء، فأمس كان
الثلاثاء وغدًا هو الخميس.
\end{example}

\section{الساعات الكاملة على الساعة}

\begin{method}[قراءة ساعة كاملة]\label{met:g1:measure:clock}
حين يشير العقرب الطويل مستقيمًا إلى الأعلى نحو $12$، تكون الساعة
كاملة: اقرأ العدد الذي يشير إليه العقرب \emph{القصير}. العقرب
القصير على $3$ والطويل على $12$: الساعة $3$.
\end{method}

\begin{omfigure}
\begin{tikzpicture}[scale=0.95]
  \draw[thick] (0,0) circle (1.3);
  \foreach \h in {1,...,12} {
    \node[font=\footnotesize] at ({90-\h*30}:1.05) {\h};
    \draw ({90-\h*30}:1.22) -- ({90-\h*30}:1.3);
  }
  \draw[very thick, omDef] (0,0) -- (0:0.62);
  \draw[thick, omThm] (0,0) -- (90:1.02);
  \node[below, font=\small] at (0,-1.55) {الساعة $3$};
  \begin{scope}[xshift=4.4cm]
  \draw[thick] (0,0) circle (1.3);
  \foreach \h in {1,...,12} {
    \node[font=\footnotesize] at ({90-\h*30}:1.05) {\h};
    \draw ({90-\h*30}:1.22) -- ({90-\h*30}:1.3);
  }
  \draw[very thick, omDef] (0,0) -- ({90-8*30}:0.62);
  \draw[thick, omThm] (0,0) -- (90:1.02);
  \node[below, font=\small] at (0,-1.55) {الساعة $8$};
  \end{scope}
\end{tikzpicture}

{\small الساعات الكاملة: العقرب الأحمر الطويل على $12$، والعقرب
الأزرق القصير يعطي الساعة. (وأنصاف الساعات وأرباعها في
\cref{ch:g2:measure}.)}
\end{omfigure}

\section{النقود الصغيرة}

\begin{example}[تكوين مبلغ]\label{ex:g1:measure:coins}
بقطع من فئة $1$ و $2$ و $5$، يمكن تكوين مبالغ كثيرة بطرق عدة.
لتكوين $7$:
\[
5 + 2, \qquad
5 + 1 + 1, \qquad
2 + 2 + 2 + 1 .
\]
والبدء بالقطع الأكبر هو الأسرع عادة، وهو يستعمل أقل عدد من
القطع.
\end{example}

\begin{omfigure}
\begin{tikzpicture}[scale=0.85]
  \draw[thick, fill=omMeth!20] (0,0) circle (0.5);
  \node[font=\small] at (0,0) {$5$};
  \node[font=\large] at (1.0,0) {$+$};
  \draw[thick, fill=omDef!20] (2.0,0) circle (0.42);
  \node[font=\small] at (2.0,0) {$2$};
  \node[font=\large] at (3.0,0) {$=$};
  \node[font=\small] at (4.0,0) {$7$};
\end{tikzpicture}

{\small قطعة من فئة خمسة وقطعة من فئة اثنين تعطيان سبعة.}
\end{omfigure}

\section{تمارين}

\begin{exercise}[$\star$]\label{exo:g1:measure:1}
صُفّ ثلاثة أقلام مطابقة من طرف واحد. أيّها الأطول؟ وأيّها
الأقصر؟
\end{exercise}

\begin{exercise}[$\star$]\label{exo:g1:measure:2}
قِس بمسطرتك: الضلع الطويل لهذا الكتاب؛ وإصبعك الصغير؛ وملعقة.
وأجب بالسنتيمترات الكاملة.
\end{exercise}

\begin{exercise}[$\star$]\label{exo:g1:measure:3}
سمِّ اليوم: بعد الاثنين؛ وقبل الأحد؛ وبعد الجمعة بيومين.
\end{exercise}

\begin{exercise}[$\star$]\label{exo:g1:measure:4}
اليوم هو الخميس. فما كان أمس؟ وما يكون غدًا؟ وما يكون بعد ثلاثة
أيام؟
\end{exercise}

\begin{exercise}[$\star$]\label{exo:g1:measure:5}
اُرسم ساعة تشير إلى $6$. أين العقرب القصير؟ وأين العقرب
الطويل؟
\end{exercise}

\begin{exercise}[$\star$]\label{exo:g1:measure:6}
اقرأ الساعة الكاملة: العقرب القصير على $5$ والطويل على $12$؛
والعقرب القصير على $11$ والطويل على $12$.
\end{exercise}

\begin{exercise}[$\star$]\label{exo:g1:measure:7}
كوِّن $8$ بقطع من فئة $1$ و $2$ و $5$، بطريقتين مختلفتين.
\end{exercise}

\begin{exercise}[$\star$]\label{exo:g1:measure:8}
عند لؤي قطعة من فئة $5$ وثلاث قطع من فئة $2$. كم من المال عنده
كله؟
\end{exercise}

\begin{exercise}[$\star\star$]\label{exo:g1:measure:9}
تكلّف الحلوى $3$. تدفع زينة بقطعة من فئة $5$. كم تسترد؟ (عُدّ
صعودًا من $3$ إلى $5$.)
\end{exercise}
